% Capítulo textual do trabalho.

\chapter{Organização do Trabalho}
\label{cap:organizacao}

Este capítulo apresenta as principais seções de um trabalho acadêmico, discute a hierarquia de seções  e fornece exemplos de como implementar a sua organização. A organização de um trabalho acadêmico é um ponto fundamental para garantir a clareza e a estrutura lógica do texto~\cite{severino}.

Normalmente, um trabalho acadêmico possui a seguinte divisão de capítulos:

\begin{alineas}
    \item Introdução: estabelece o contexto do trabalho e apresenta o problema de pesquisa;
    \item Fundamentação teórica: explora os conceitos e teorias que embasam a pesquisa, proporcionando uma base teórica para a compreensão do tema;
    \item Revisão da Literatura: analisa e discute estudos anteriores relacionados ao tema, identificando lacunas que a pesquisa busca preencher;
    \item Metodologia: descreve os métodos e técnicas utilizados para conduzir a pesquisa, permitindo a reprodução do estudo por outros pesquisadores;
    \item Resultados Experimentais: apresenta e analisa os dados obtidos durante a pesquisa, evidenciando os achados e sua relevância para o tema investigado;
    \item Conclusão: resume os principais resultados, discute suas implicações, e sugere direções para futuras pesquisas ou aplicações práticas.
\end{alineas}

A organização adequada de um trabalho acadêmico não apenas melhora a legibilidade, mas também facilita a navegação pelo documento. O uso correto dos comandos de seção, subseção e subsubseção ajuda a estruturar o conteúdo de forma lógica e coesa, orientando o leitor através dos diferentes tópicos abordados. Além disso, a numeração automática das seções e subseções torna o trabalho mais profissional e alinhado às normas acadêmicas~\cite{gil}.
